\documentclass[../../main.tex]{subfiles}
\begin{document}

\tikzstyle{State}=[circle, draw, thick]

\begin{problem}
\begin{itemize}
    \item [(1)] 利用 prim 算法，求下面无向图的最小生成树。
    \item [(2)] 利用 Kruskal 算法，求下面无向图的最小生成树。
    \item [(3)] 利用破圈法，求下面无向图的最小生成树。
\end{itemize}
\end{problem}
\begin{center}
\begin{tikzpicture}
% Nodes
\node[State] (A) at ( 0  , 3) {A};
\node[State] (B) at (-1.5, 2) {B};
\node[State] (C) at ( 0  , 1) {C};
\node[State] (D) at ( 1.5, 2) {D};
\node[State] (E) at (-1.5, 0) {E};
\node[State] (F) at ( 1.5, 0) {F};
% Edges with weights
\draw[thick] (A) -- (B) node[midway, above left]  {6};
\draw[thick] (A) -- (C) node[midway, right]       {1};
\draw[thick] (A) -- (D) node[midway, above right] {5};
\draw[thick] (B) -- (C) node[midway, above right] {5};
\draw[thick] (B) -- (E) node[midway, left]        {3};
\draw[thick] (C) -- (E) node[midway, above left]  {6};
\draw[thick] (C) -- (D) node[midway, above left]  {5};
\draw[thick] (C) -- (F) node[midway, above right] {4};
\draw[thick] (D) -- (F) node[midway, right]       {2}; 
\draw[thick] (E) -- (F) node[midway, below]       {6};
\end{tikzpicture}
\end{center}

\begin{solution}
\begin{itemize}
    \item [(1)] 利用 prim 算法，求下面无向图的最小生成树。
    
    从 A 开始选择: A-C, C-F, D-F, C-B, B-E
    \begin{center}
    \begin{tikzpicture}
    % Nodes
    \node[State] (A) at ( 0  , 3) {A};
    \node[State] (B) at (-1.5, 2) {B};
    \node[State] (C) at ( 0  , 1) {C};
    \node[State] (D) at ( 1.5, 2) {D};
    \node[State] (E) at (-1.5, 0) {E};
    \node[State] (F) at ( 1.5, 0) {F};
    % Edges with weights
    \draw[thick] (A) -- (C) node[midway, right]       {1};
    \draw[thick] (C) -- (F) node[midway, above right] {4};
    \draw[thick] (D) -- (F) node[midway, right]       {2}; 
    \draw[thick] (B) -- (C) node[midway, above right] {5};
    \draw[thick] (B) -- (E) node[midway, left]        {3};
    \end{tikzpicture}
    \end{center}
    \item [(2)] 利用 Kruskal 算法，求下面无向图的最小生成树。
    
    从小到大选择: A-C, D-F, B-E, C-F, B-C
    \begin{center}
    \begin{tikzpicture}
    % Nodes
    \node[State] (A) at ( 0  , 3) {A};
    \node[State] (B) at (-1.5, 2) {B};
    \node[State] (C) at ( 0  , 1) {C};
    \node[State] (D) at ( 1.5, 2) {D};
    \node[State] (E) at (-1.5, 0) {E};
    \node[State] (F) at ( 1.5, 0) {F};
    % Edges with weights
    \draw[thick] (A) -- (C) node[midway, right]       {1};
    \draw[thick] (D) -- (F) node[midway, right]       {2}; 
    \draw[thick] (B) -- (E) node[midway, left]        {3};
    \draw[thick] (C) -- (F) node[midway, above right] {4};
    \draw[thick] (B) -- (C) node[midway, above right] {5};
    \end{tikzpicture}
    \end{center}
    \item [(3)] 利用破圈法，求下面无向图的最小生成树。
    
    按照顺序将以下边移除: A-D, A-B, C-E, E-F, C-D
    \begin{center}
    \begin{tikzpicture}
    % Nodes
    \node[State] (A) at ( 0  , 3) {A};
    \node[State] (B) at (-1.5, 2) {B};
    \node[State] (C) at ( 0  , 1) {C};
    \node[State] (D) at ( 1.5, 2) {D};
    \node[State] (E) at (-1.5, 0) {E};
    \node[State] (F) at ( 1.5, 0) {F};
    % Edges with weights
    \draw[thick] (A) -- (C) node[midway, right]       {1};
    \draw[thick] (B) -- (C) node[midway, above right] {5};
    \draw[thick] (B) -- (E) node[midway, left]        {3};
    \draw[thick] (C) -- (F) node[midway, above right] {4};
    \draw[thick] (D) -- (F) node[midway, right]       {2}; 
    \end{tikzpicture}
    \end{center}
\end{itemize}
\end{solution}

\end{document}